\section{State of the art}\label{sec:sota}

\subsection{UV-induced defects in silica and their relevance at 369.5\,nm}\label{sec:sota-bulk}
\paragraph{Energetics} A \SI{369.5}{nm} photon carries \SI{3.36}{eV}; two photons carry \SI{6.71}{eV}.
The absorption edge of amorphous \ce{SiO2} lies near \SI{8}{eV}~\cite{Kajihara2010}, so neither one nor
two photons can excite the perfect glass across its gap. Excimer-laser studies show that E$'$ centres form
in a two-photon process at \SI{5.0}{eV} and \SI{6.4}{eV}~\cite{Arai1988,Tsai1988}. XeCl light
(\SI{4.0}{eV}, two-photon energy \SI{8.0}{eV}) created no centres in dry synthetic
silica~\cite{Arai1988}. At \SI{369.5}{nm} only processes that start at a pre-existing defect, an impurity
or a surface site are therefore expected: one-photon absorption into a weak band, or stepwise two-photon
absorption through a long-lived intermediate state.

\paragraph{The relevant defects} \Cref{tab:defects} lists the main point defects of silica
glass~\cite{Skuja1998,Griscom1991,Skuja2005,Skuja2012,Girard2019}. Most absorb well above
\SI{3.36}{eV}, but several channels exist at \SI{369.5}{nm}:
\begin{itemize}
  \item the spin-forbidden singlet--triplet band of the twofold-coordinated silicon centre ODC(II) at
        \SI{3.15}{eV}, whose long-lived (\si{ms}) triplet state can absorb a second photon;
  \item the \SI{6.9}{eV} band of ODC(II) and the onset of the Schumann--Runge absorption of interstitial
        \Oxy{} (shifted to $\approx\SI{6.5}{eV}$ in the glass), both reachable by two photons;
        photolysed interstitial \Oxy{} yields mobile oxygen atoms, interstitial ozone and peroxy
        radicals~\cite{Awazu1990,Kajihara2004a,Kajihara2005a,Skuja2000};
  \item chlorine species left by the manufacturing process: interstitial \ce{Cl2} absorbs at
        \SI{3.78}{eV} with a width of \SI{0.6}{eV}, and atomic \ce{Cl^0} at \SI{3.26}{eV}. In pure-silica-core
        fibres, much of the radiation-induced loss above \SI{2.5}{eV} is chlorine-related~\cite{Skuja2017,Girard2019}.
\end{itemize}
In pure-silica-core single-mode fibre under \SI{405}{nm} CW light at \si{MW/cm^2}, bleaching of ODC(II) and
growth of E$'$ centres have indeed been observed~\cite{Gonschior2013JNCS}. Germanium-related centres
would add a one-photon channel through the GeODC triplet band near \SI{3.8}{eV}. This is the origin of
the photosensitivity of Ge-doped fibres at near-UV and, via two photons, visible
wavelengths~\cite{Hill1978,Starodubov1997,Grubsky1999}. The group's PM-S350-HP fibre has a pure-silica core, so these
channels should be absent.

\begin{table}[t]
\centering\small
\caption{Main point defects in silica glass: absorption and photoluminescence (PL) bands, in eV, after
\cite{Skuja1998,Skuja2011,Skuja2017,Girard2019}. Values in bold are within reach of one (\SI{3.36}{eV})
or two (\SI{6.71}{eV}) photons at \SI{369.5}{nm}.}
\label{tab:defects}
\begin{tabularx}{\textwidth}{@{}L{4.6cm} L{3.6cm} L{2.4cm} Y@{}}
\toprule
Defect & Absorption (eV) & PL (eV) & Comment \\
\midrule
E$'$ centre, $\equiv$Si$^{\bullet}$ & 5.8 & -- & dominant UV-induced centre \\
NBOHC, $\equiv$Si--O$^{\bullet}$ & 1.97, 4.8 (continuum 4--7.8) & 1.91 & \SI{3.36}{eV} lies in its low-absorption gap \\
ODC(I), $\equiv$Si--Si$\equiv$ & 7.6 & 4.4 & precursor of E$'$ \\
ODC(II), $=$Si: & 5.05, \textbf{6.9}; triplet \textbf{3.15} & 4.4, 2.75 & triplet band and long-lived T$_1$ state \\
peroxy radical, $\equiv$Si--O--O$^{\bullet}$ & $\approx$2.0, 4.8 & -- & forms from NBOHC + O, or E$'$ + \Oxy{} \\
interstitial \Oxy{} & onset \textbf{$\approx$6.5} & 0.975 & source of O atoms and \ce{O3} \\
interstitial \ce{O3} & 4.8 & -- & bleached by UV \\
interstitial \ce{Cl2}, \ce{Cl^0} & \textbf{3.78} (FWHM 0.6), \textbf{3.26} & -- & chlorine from manufacture \\
GeODC (Ge-doped cores only) & 5.15; triplet \textbf{$\approx$3.8} & 4.3, 3.2 & Ge photosensitivity \\
\bottomrule
\end{tabularx}
\end{table}

\paragraph{Hydrogen, OH and fibre manufacture} Molecular hydrogen passivates E$'$ centres and NBOHCs. The
E$'$ + \Htwo{} reaction has an activation energy of only \SI{0.38}{eV} and proceeds at the same rate for
surface E$'$ centres~\cite{Radzig1995,Messina2007}. Under ArF and \ce{F2} irradiation \Htwo{} can also
promote ODC(I) formation~\cite{Ikuta2002}. Fibre drawing creates E$'$ centres, NBOHCs (the \SI{630}{nm}
band, which grows with draw tension) and peroxy radicals~\cite{Hanafusa1985,Hibino1986,Hibino1987}.
Strained Si--O--Si bonds are the main intrinsic precursors~\cite{Kajihara2008b}. Decades of work on
multimode UV fibres have led from hydrogen-loaded to hydrogen-free, high-OH designs with improved deep-UV
stability~\cite{Karlitschek1995,Karlitschek1998a,Klein2001,Khalilov2006,KleinKhalilov2015}. Loading with
\Htwo{} and UV curing has transferred this to single-mode PCF~\cite{Colombe2014}.

\paragraph{Kinetics} Colour-centre growth in fibres typically follows a stretched exponential or a
power law in dose~\cite{Griscom1993,Griscom2001}. Two-photon creation shows up as a quadratic
dependence on intensity~\cite{Arai1988,Tsai1988}. Comparable methods for measuring and fitting
photodarkening have been standardised for Yb-doped fibres~\cite{Koponen2006,Koponen2008,Jetschke2007}.
Measuring the initial rate at three or more powers therefore gives a cheap first indication of the
mechanism.

\paragraph{Gas transport into the glass} At room temperature \Htwo{} diffuses through a
\SI{125}{\micro m} fibre in a few days~\cite{Lemaire1991,Shang2009}. Oxygen molecules move about ten
orders of magnitude more slowly. Extrapolated high-temperature data~\cite{Norton1961,Kajihara2010} give
penetration depths of 0.1--5\,nm per year (\cref{app:estimates}). The bulk of the core is therefore
isolated from the atmosphere, except for deliberately added hydrogen. Oxygen can act only at the
surface, where it dissolves readily~\cite{Kajihara2004c}.

\paragraph{Reactive surfaces} Freshly fractured or ground silica carries reactive sites: silicon and
oxygen dangling bonds, silylene groups and strained two-membered rings~\cite{Radzig1993,Radzig1998,Radzig2000,Ceresoli2000}.
These react with \water{} (which opens the strained rings~\cite{Rimola2008}), with \Htwo{} and with
\Oxy{}. Polished end faces are such mechanically activated surfaces. Oxygen, water and hydrogen
therefore determine what UV-generated surface defects turn into: peroxy radicals, silanols or Si--H.
These products absorb mainly above \SI{4}{eV}, not at \SI{3.36}{eV}. Spectroscopically, H1 is therefore
not the most likely mechanism, which is one more reason to test it directly rather than assume it.


\subsection{What has been observed in single-mode fibres carrying UV and blue light}\label{sec:sota-reports}
\Cref{tab:reports} collects the most relevant reports. Two kinds of degradation emerge clearly.
\emph{Bulk solarisation} dominates in the deep UV. It is irreversible, it can set in at microwatt powers
in fused-silica fibres~\cite{Colombe2014}, and what light remains is carried in a speckled mode
pattern~\cite{Marciniak2017}. Below about \SI{300}{nm}, unprotected photonic-crystal fibres lose most of
their transmission within hours at milliwatt powers~\cite{Colombe2014,Gebert2014}. \emph{End-face
degradation} is reversible by re-cleaving or re-polishing, and it is the more relevant kind at near-UV and
blue wavelengths. It takes at least three forms:
\begin{itemize}
  \item \textbf{Deposits.} Colombe et al.\ mention UV-induced carbon deposition on input
        faces~\cite{Colombe2014}. Marciniak et al.\ saw deposits on large-mode-area PCF faces at
        \SI{313}{nm} with input powers as low as \SI{10}{mW}. Polishing removed them, and a slight
        overpressure of dry \Ntwo{} at the collimators prevented them for months~\cite{Marciniak2017}. They
        attributed the deposits to particles and oil drawn into the beam by the ``optical-tweezer effect''
        and then burned in. Purged and unpurged ends of the same fibre behaved differently. A Duke group
        reports the same use of nitrogen purging against particle deposition at UV
        wavelengths~\cite{YuJ2026}.
  \item \textbf{Restructuring of the silica.} In the most detailed study of CW blue light in small-core
        single-mode fibres (\SI{405}{nm}, mode-field diameter \SIrange{2}{3}{\micro m}), Gonschior et al.\
        found no deposits. Instead, laser-induced periodic surface structures formed on both faces:
        protrusions of \SIrange{0.5}{0.7}{\micro m} after 9--11 days, ripples of \SIrange{250}{400}{nm}
        period aligned with the polarisation, and a porous, lower-density silica layer. EDX and ToF-SIMS
        found no foreign elements~\cite{Gonschior2013JNCS,Gonschior2013PhD}. The loss grew faster than
        linearly in time, typically \SIrange{0.2}{1}{dB/day} at \SIrange{74}{95}{mW}. The output beam
        quality degraded ($M^2$ from 1.2 to 5 in 12 days). Cleaved faces degraded four times more slowly
        than polished ones~\cite{Gonschior2013PhD,Gonschior2014OE}. The structures appeared at
        \SIrange{0.5}{2.5}{MW/cm^2}, three orders of magnitude below the optical damage
        threshold~\cite{Gonschior2013JNCS}.
  \item \textbf{Slow drifts of unclear origin.} At \SI{369}{nm}, Sandia found lensed-fibre-to-chip coupling stable at
        \SIrange{7}{24}{mW} on short time scales but degrading over hours, with surface contamination of
        the fibre or chip facets among the suspected causes~\cite{Kwon2024}.
\end{itemize}
Two further observations bear directly on the oxygen hypothesis. Gonschior ran a PM single-mode fibre for
50 days at \SI{250}{mW} (\SI{405}{nm}) between \emph{nitrogen-purged, sealed} collimators. The fibre
lost only about \SI{1}{dB}, but cone-like structures still formed on both faces~\cite{Gonschior2013PhD}.
Inert gas alone therefore did not stop the restructuring at that (much higher) intensity. Second, in
ion-trap vacuum systems, \SIrange{369}{370}{nm} light caused losses in oxide mirror coatings that were
reversed by admitting oxygen or air~\cite{Gangloff2015,Ballance2017}. Absence of oxygen can therefore
\emph{harm} UV-illuminated oxide surfaces.

\begin{table}[t]
\centering\footnotesize
\caption{Selected reports on UV/blue degradation of single-mode fibres and related surfaces. PCF:
photonic-crystal fibre; LMA: large mode area; HC: hollow core.}
\label{tab:reports}
\begin{tabularx}{\textwidth}{@{}L{1.9cm} L{1.2cm} L{2.6cm} Y L{3.6cm}@{}}
\toprule
Ref. & $\lambda$ (nm) & Fibre / surface & Observation & Mitigation that worked \\
\midrule
\cite{Colombe2014} & 313, 280 & LMA PCF & bulk solarisation within hours at \si{mW}; carbon on input face
  & \Htwo{} loading + UV curing; collapsed, angled ends \\
\cite{Gebert2014} & 280 & solid-core PCF; kagom\'e HC & solid core $-$80\,\% in 3\,h at \SI{2}{mW}; HC unchanged after \SI{100}{h} at \SI{15}{mW}
  & hollow core \\
\cite{Yamamoto2009} & 250 & LMA PCF & damage mainly at the input; partly reversed by re-cleaving & low power ($<\SI{0.3}{mW}$) \\
\cite{Gonschior2013JNCS,Gonschior2013PhD} & 405 & SM and PM, MFD 2--3\,\textmu m & LIPSS and porous silica on both faces; $M^2$ grows; no foreign elements
  & end caps or launch fibre: $<\SI{0.03}{dB/day}$ \\
\cite{Marciniak2017} & 313 (369) & LMA-PM PCF & particle/oil deposits from \SI{10}{mW}; removed by polishing & dry \Ntwo{} overpressure \\
\cite{Kwon2024} & 369 & lensed fibre to chip & slow (hour-scale) coupling loss & -- \\
\cite{Gangloff2015,Ballance2017} & 369--422 & oxide mirror coatings in vacuum & UV-induced loss in vacuum & reversed by \Oxy{} or air \\
\cite{Hippler2016,Kokkinos2017} & 355 & optics with organic contaminants & deposit growth in vacuum & \Oxy{} partial pressure \\
\bottomrule
\end{tabularx}
\end{table}

\begin{keybox}[What is not known]
We have found no published measurement of the degradation of step-index, pure-silica-core PM fibre at
\SI{369.5}{nm} and a few milliwatts, the configuration used in the group. No study in the trapped-ion
community has treated \Oxy{} as a controlled variable. Nitrogen purging is used to keep particles away,
not to remove oxygen. This is the gap the project fills.
\end{keybox}

\subsection{Fibre technologies for near-UV delivery}\label{sec:sota-fibres}
\paragraph{The group's fibre} The cables known in the group as ``375PM'' are Thorlabs P3-375PM-FC
patch cables (\SIrange{1}{10}{m}), FC/APC at both ends, aligned at \SI{375}{nm}, with no AR-coated
option~\cite{ThorlabsP3375PM}. They are built on PM-S350-HP: a PANDA fibre with a \emph{pure-silica
core}, mode-field diameter \SI{2.3}{\micro m} at \SI{350}{nm} (about \SI{2.4}{\micro m} at
\SI{369.5}{nm}), NA~0.12 and second-mode cut-off $\leq\SI{340}{nm}$~\cite{ThorlabsPMS350}. A
germanium-free core avoids the strong photosensitivity of Ge-doped silica, and Thorlabs advertises
pure-silica cores for their resistance to photodarkening. The same catalogue nevertheless states that all
fibres eventually photodarken in the UV and should be treated as consumables. For CW light at an
air--glass face of undoped silica, Thorlabs quotes a practical safe intensity of about
\SI{250}{kW/cm^2}~\cite{ThorlabsDamage}. With Thorlabs' convention this corresponds to about \SI{11}{mW}
for a \SI{2.4}{\micro m} mode. The group's few milliwatts are below that, yet the faces still degrade,
which illustrates that ``damage threshold'' figures say little about slow photochemical processes.
Thorlabs no longer sells the bare fibre; Coherent (Nufern) offers a fibre with the same part number. The
\SI{405}{nm} PM cables use PM-S405-XP, whose cut-off ($380\pm\SI{20}{nm}$) may allow a second mode at
\SI{369.5}{nm}. They are therefore no substitute.

\paragraph{Large-mode-area photonic-crystal fibre with hydrogen loading} Endlessly single-mode PCFs
spread the light over roughly ten times the mode area. Loading them with \Htwo{} (about \SI{10}{MPa},
room temperature, 4--6 days) and curing them with UV light before the hydrogen escapes makes them resistant
to deep-UV solarisation for hundreds of hours at \SI{125}{mW} (\SI{313}{nm})~\cite{Colombe2014,NISTrecipe}.
Hydrogen loading alone does not confer resistance unless curing occurs before the hydrogen has diffused
out~\cite{Colombe2014}. Collapsing the holes over a few hundred micrometres and angle-cleaving the ends
protects the faces and suppresses back-reflections. PM versions reach polarisation extinction ratios of
30--32\,dB at \SI{369}{nm} after curing~\cite{Marciniak2017}. These fibres are not off-the-shelf
patch cords, and their faces still collect deposits without a purge~\cite{Marciniak2017}.

\paragraph{Hollow-core fibres} Kagom\'e and antiresonant hollow-core fibres guide almost all the light in
gas, so the glass hardly sees it. They have transmitted \SI{280}{nm} CW light without change for
\SI{100}{h}~\cite{Gebert2014} and \SI{320}{nm} light at up to \SI{60}{mW} without degradation, at about
\SI{0.13}{dB/m}~\cite{Dorer2024}. Losses below \SI{10}{dB/km} have been reported at
\SI{369}{nm}~\cite{Osorio2023}, and guidance extends down to \SI{190}{nm}~\cite{Mears2024}. Caveats:
there is no PM hollow-core fibre, ammonium chloride can grow on cleaved hollow-core end
faces~\cite{Rikimi2020}, and reversible photodarkening has been reported in UV-guiding hollow-core
fibres~\cite{Mears2024SPIE}.

\paragraph{End caps} A coreless silica cap spliced to the fibre lets the beam expand before it reaches the
glass--gas interface. At \SI{369.5}{nm} a \SI{300}{\micro m} cap reduces the face intensity about
280-fold (\cref{app:estimates}). At \SI{405}{nm}, a short launch fibre and end caps reduced the
damage rate from \SIrange{0.2}{1}{dB/day} to below \SI{0.03}{dB/day}~\cite{Gonschior2013PhD}. End caps
attack the intensity-driven surface processes at their root. Together with the atmosphere control
studied here they are the most promising engineering solution, and the obvious next step if the
atmosphere alone is not enough.

\subsection{Surface processes and the role of the atmosphere}\label{sec:sota-surface}
Four communities have studied UV-driven degradation of optical surfaces in controlled atmospheres:
space lidar, UV frequency-conversion cavities, semiconductor lasers and laser modules, and
high-power fibre connectors. Their results are the best available guide to what the atmosphere around a
fibre end face will do.

\paragraph{Laser-induced contamination in space lasers} A prototype of the \SI{355}{nm} ALADIN laser for
ESA's Aeolus mission lost \SI{50}{\percent} of its transmission in less than six hours. Outgassed organics
were carbonised on the optics by the UV light, and even a few parts per billion of organics caused
contamination~\cite{ESA2018Aeolus}. The flight laser therefore operates with a low-pressure oxygen purge
of a few tens of pascals~\cite{Wernham2010,ESA2018Aeolus}. Systematic tests at DLR showed that an
\Oxy-to-contaminant pressure ratio above 50:1 drastically reduces deposit growth, above 1000:1 no
transmission loss was detectable, and existing deposits were completely removed by irradiation in
several tens of pascals of oxygen~\cite{Hippler2016,Tighe2008,Riede2011,Schroeder2013}. The ESA/DLR
guidance for ground operation is explicit: ``operation with normal dry air is to be preferred to operation
in nitrogen''~\cite{Mondin2017}. The same paper notes that silicones are not removed by oxidation, that CW
irradiation can produce more deposit than pulsed irradiation, and that high-purity bottled gas alone does
not prevent absorbing ``black dots''. In sealed \SI{1064}{nm} lasers, trace contamination from
ultra-high-purity nitrogen cylinders was enough to cause damage above about
\SI{60}{\celsius}~\cite{Hovis1994}. With \SI{300}{ppm} toluene
present, optics survived $8\times10^3$ pulses in \Ntwo, $1.8\times10^4$ in synthetic air, and more than
$10^6$ in air with a few thousand ppm of water or alcohol~\cite{Weiller2010}.

\paragraph{CW evidence at our intensity} Parthasarathy et al.\ studied contamination at a diamond
surface in vacuum under \SI{1.5}{mW} of CW \SI{532}{nm} light focused to about \SI{220}{kW/cm^2}, an
intensity close to that on the group's fibre faces~\cite{Parthasarathy2024}. Deposits grew
\emph{fastest} at an oxygen partial pressure of about \SI{5e-3}{mbar}. Between \SI{2e-3}{mbar} and
\SI{2}{mbar} their onset was delayed, and above \SI{2}{mbar} none formed. Humidity slowed the growth. At
atmospheric pressure, \SIrange{1}{1000}{ppm} of \Oxy{} corresponds to \SIrange{0.001}{1}{mbar}. An
``almost oxygen-free'' cell is therefore exactly in the regime where contamination was promoted, whereas
air is a hundred times above the inhibition threshold. How these vacuum results transfer to
\SI{1}{atm} of nitrogen is not known. That is one of the questions the \arm{O2-x} arms answer.

\paragraph{UV conversion cavities} Groups producing CW UV light for trapped ions routinely flush their
nonlinear crystals with oxygen. At NIST the \SI{313}{nm} BBO crystal sits in an enclosure ``purged with a
gentle flow of dry oxygen'', there mainly against moisture-induced fogging~\cite{Wilson2011}. At PTB oxygen is led directly onto the crystal facets ``to
prevent light-induced damage''~\cite{Hannig2018}. At \SI{243}{nm}, cavity mirrors degraded within minutes
below about \SI{0.7}{mbar} of oxygen, and the deposit could be removed with organic
solvents~\cite{Cooper2018,Burkley2021}. A recent watt-level \SI{230}{nm} source uses filtered dry air
because ``pure nitrogen removes this cleaning pathway''~\cite{Cai2026}. One industrial patent reports the
opposite for BBO: clean dry air alone degraded the crystal within 100 hours, and adding about \SI{1000}{ppm}
of water to the purge cured this~\cite{KLApatent}.

\paragraph{Semiconductor lasers and fibre-coupled laser modules} GaN laser diodes operated in pure
nitrogen grow a thick, highly absorbing deposit on the facet, which oxygen removes~\cite{Schoedl2005}.
Carbon deposition~\cite{Kim2008} and the growth of \ce{SiO_x} layers~\cite{Freier2024,WangXW2021} are
documented facet failure modes. Package patents mostly specify oxygen-containing fills: more than
\SI{100}{ppm}~\cite{Corningpatent}, at least \SI{5}{\percent}~\cite{Sanyopatent}, or
\SIrange{1}{18}{\percent}~\cite{Sharppatent}. For \SI{405}{nm} lasers, oxygen is said to react with
photo-decomposed silicon compounds before they deposit on the facet~\cite{Panasonicpatent}. The important
exception concerns \emph{fibre} end faces. In a fibre-coupled \SI{410}{nm} module, reliability
\emph{deteriorated} at oxygen concentrations of \SI{100}{ppm} or more, because silicon-based compounds
deposited at the fibre end face. Silicone parts were the principal source, and the patent specifies
\SIrange{1}{100}{ppm} \Oxy~\cite{Fujifilmpatent2006}. A related patent on siloxane-derived \ce{SiO_x}
deposits at a fibre input recommends silicone-free materials and air admitted only through activated
carbon~\cite{Ricohpatent}. Another describes a ``dust collection effect'' where the optical density is
high~\cite{Fujifilmpatent2007}.

\paragraph{Can 369.5\,nm light make ozone?} Splitting \Oxy{} needs photons below \SI{242}{nm} (the
Herzberg continuum~\cite{Yoshino1988}). UV/ozone cleaning therefore uses the \SI{185}{nm} and
\SI{254}{nm} mercury lines~\cite{Vig1985}. Two \SI{369.5}{nm} photons carry enough energy, but no
two-photon cross-section near \SI{370}{nm} is known, and any effect would be confined to the
\SI{\approx20}{\micro m} next to the face. Oxygen suppresses contamination even where it cannot be
photodissociated~\cite{Parthasarathy2024,Hippler2016}. The relevant chemistry is light-assisted oxidation
at the surface, not gas-phase ozone.

\paragraph{Particles and high-power connectors} Dielectric particles are trapped by the optical gradient
force in air~\cite{Ashkin1986,Omori1997}, and absorbing particles are moved by photophoretic
forces~\cite{Shvedov2009}. Fibre connectors carrying watts at \SI{1550}{nm} fail mainly because of
absorbing contamination on the face, at intensities comparable to ours~\cite{DeRosa2002,Seo2003,Domingues2011}.
The \SI{313}{nm} fibre deposits of Marciniak et al.\ were attributed to this mechanism~\cite{Marciniak2017}.

\begin{keybox}[Oxygen: friend or foe?]
\begin{itemize}[leftmargin=1em]
  \item \textbf{Hydrocarbon contamination:} every community finds oxygen (and often water vapour)
        \emph{protective}. Nitrogen and vacuum favour absorbing carbon deposits, most strongly at
        trace \Oxy{} levels.
  \item \textbf{Siloxane contamination:} oxidation turns silicones into non-volatile \ce{SiO_x}. At a
        \SI{410}{nm} fibre end face, oxygen above \SI{100}{ppm} \emph{accelerated} such deposits. This is
        the strongest support in the literature for the group's hypothesis.
  \item \textbf{Particles and silica restructuring:} not expected to depend on oxygen. A clean, slightly
        pressurised gas keeps particles away.
\end{itemize}
The outcome for the group's fibres therefore depends on which contaminants are present, and finding
that out is the first task (Phase~0). A sealed, nearly oxygen-free cell that contains outgassing
polymers is a high-risk configuration. The experiment tests it against clean, oxygen-containing gas
rather than assuming it helps.
\end{keybox}


\subsection{Creating and verifying a clean, oxygen-free atmosphere}\label{sec:sota-atmosphere}
\paragraph{Oxygen absorbers} Iron-based absorbers (Mitsubishi Gas Chemical AGELESS) reduce the oxygen
in a sealed barrier container to $\leq\SI{0.1}{\percent}$. The container wall should let through no more
than \SI{20}{mL} of oxygen per square metre, day and atmosphere of pressure difference~\cite{MGCAgeless}. The number in the type name is the volume of
\Oxy{} absorbed: a ZP-50 binds \SI{50}{mL}. Some types need moisture from their surroundings. The
self-reacting Z types work best above about \SI{55}{\percent} RH and release humidity into the
container~\cite{Daniel1993}. The reaction is exothermic, and the sachets lose capacity within minutes in
air. Lambert et al.\ give a useful rule for a leaking box: the equilibrium concentration is
$[\Oxy]\approx L/(12.7\,C)$, where $L$ is the ingress in ppm per day and $C$ is the absorber capacity
divided by the box volume~\cite{Lambert1992}. The RP System of the same manufacturer exists in two
types. RP-A absorbs oxygen, water and corrosive gases and drives the humidity to nearly zero. RP-K
leaves the humidity unchanged but can release organic acids if overloaded~\cite{MGCRP}. Both reach
$\leq\SI{0.1}{\percent}$ \Oxy{} within a day and shed fine powder, so the sachet belongs behind a
membrane. Colour indicators change near \SIrange{0.1}{0.5}{\percent} \Oxy{} and show only gross
failures. Vacuum non-evaporable getters need activation at \SIrange{350}{450}{\celsius} and would
saturate almost immediately in a cell at atmospheric pressure~\cite{SAESSorbAC}. In heated form they serve
as in-line gas purifiers. Commercial point-of-use purifiers deliver inert gas with
$<\SI{100}{ppt}$ \Oxy{} and \water{}~\cite{EntegrisHX}. A palladium catalyst with a few per cent \Htwo{}
holds anaerobic chambers at \SIrange{0}{5}{ppm} \Oxy~\cite{Coy}, an option if forming gas is used anyway.

\paragraph{Measuring oxygen, water and organic vapour} Proving $<\SI{100}{ppm}$ needs an instrument that
resolves ppm.
\begin{itemize}
  \item Ultra-trace luminescent sensor spots read through a window: e.g.\ PreSens SP-PSt9,
        \SIrange{0}{200}{ppm}, detection limit \SI{0.5}{ppm}~\cite{PreSensPSt9}. They are ideal for
        sealed cells. The common trace spots with a \SI{50}{ppm} detection limit are only marginal for this
        purpose~\cite{PyroTROX}, and percent-level sensors are useless here~\cite{LuminOx}.
  \item Galvanic trace analysers (detection limit \SI{\sim0.05}{ppm}) need a sample flow of
        \SIrange{0.5}{1}{L/min}. They suit the exhaust of purged arms rather than sealed
        cells~\cite{GPR1200}.
  \item Zirconia sensors read wrongly in hydrogen-containing gas and are poisoned by silicone
        vapour~\cite{PSTzirconia}.
\end{itemize}
Luminescent spots must be shielded from the UV beam. Capacitive humidity sensors resolve trends down to
about \SI{1}{\percent} RH (\SI{\approx230}{ppmv})~\cite{SHT4x}. Low-cost metal-oxide VOC sensors depend on
oxygen adsorbed on the oxide~\cite{Barsan2001}, so they work only in oxygen-containing streams. Organic
vapour in oxygen-free streams is best sampled onto sorbent tubes and analysed by TD-GC/MS.

\paragraph{Gas purity} Nitrogen 5.0 contains up to \SI{3}{ppm} \Oxy, \SI{5}{ppm} \water{} and
\SI{0.2}{ppm} hydrocarbons; 6.0 grade one sixth of that~\cite{LindeN2}. ``Zero grade'' air may contain
\SI{5}{ppm} of hydrocarbons~\cite{BOCzero}, whereas ultra-zero grades are specified at
\SI{0.1}{ppm}~\cite{AirgasUZ}. The hydrocarbon content must be specified explicitly, because
parts-per-billion levels already matter (\cref{app:estimates}).

\paragraph{3D-printed enclosures} FDM parts become leak-tight with slight over-extrusion (extrusion
multiplier 0.98 to 1.2 depending on the printer), walls of at least \SI{1.1}{mm} and randomised
seams~\cite{Gordeev2018}. Thin layers matter more than infill~\cite{Leeder2025}, and laser-sintered PA12
permeates like bulk material~\cite{Liebrich2021}. Printed parts are used in vacuum systems: SLS PA12 was
characterised by Povilus et al.~\cite{Povilus2014}; FDM PEEK, SLA resins and other printed plastics have
been tested since~\cite{Stallwood2018,Radic2025,Domingues2025}. Printed metal chambers reach ultra-high
vacuum in portable quantum technology~\cite{Cooper2021}. \Cref{tab:materials} lists permeabilities and
outgassing data. Three points stand out:
\begin{itemize}
  \item Dense polymer walls let in 50--500 times more water molecules than oxygen molecules. Printed ABS
        stores \SI{0.35}{wt\percent} of gas, mostly water~\cite{Sefa2016}. A sealed ``dry'' cell
        therefore becomes wet within hours unless it contains a desiccant (\cref{app:estimates}).
  \item Standard SLA resins lose \SIrange{1.5}{2.9}{\percent} of their mass in vacuum, whereas PETG,
        PEEK, ULTEM and Antero lose about \SI{0.3}{\percent}~\cite{NASAoutgassing}. FDM printing of
        ABS releases styrene, and cured resins still emit acrylates and
        aldehydes~\cite{Azimi2016,Wojnowski2022,Pham2023}.
  \item O-rings need a surface finish of about \SI{0.4}{\micro m}, squeeze of at most \SI{30}{\percent}
        and \SIrange{60}{85}{\percent} gland fill~\cite{ParkerORing}. FDM surfaces are much rougher,
        so the seal should sit on an ironed top skin, a machined land or an SLA part. Silicone O-rings
        let in about 200 times more oxygen than FKM or butyl, and they release siloxanes.
\end{itemize}

\begin{table}[t]
\centering\small
\caption{Permeability at about \SI{23}{\celsius} (1\,Barrer $=\SI{3.35e-16}{mol.m.m^{-2}.s^{-1}.Pa^{-1}}$)
and outgassing (ASTM E595 total mass loss TML and collected volatile condensable material CVCM) of
candidate cell materials. Values compiled from \cite{Marano2022,Levitsky2019,Sonstevold2023,Liebrich2021,Peacock1980,Chiggiato2020,NASAoutgassing};
-- not found. Permeabilities of PETG, ASA, ABS and PEI should be measured by rate of rise.}
\label{tab:materials}
\footnotesize
\begin{tabularx}{\textwidth}{@{}l c c c Y@{}}
\toprule
Material & \Oxy{} (Barrer) & \water{} (Barrer) & TML / CVCM (\%) & Remark \\
\midrule
PET / PETG          & 0.034 (PET) & 130 (PET) & 0.35 / 0.02 (PETG) & good barrier; low VOC emission when printed \\
PLA                 & 0.15--0.43  & 2000--3300 & 0.70--0.75 / 0.01 & very water-permeable \\
ABS                 & --          & $\approx$2600 & 0.59--0.94 / 0.01--0.04 & emits styrene; avoid \\
PC                  & 1.0         & --         & 0.14--0.17 / 0.00 & \\
PP                  & 1.7         & --         & -- & poor \Oxy{} barrier \\
PA12 (SLS)          & 0.13--0.20  & --         & 0.75 / 0.02 & hygroscopic \\
PEEK                & 0.14        & $\approx$140 & 0.26--0.30 / $\leq$0.02 & needs a high-temp.\ printer \\
ULTEM 9085 / Antero 800NA & --    & --         & 0.33--0.37 / $\leq$0.02; 0.35 / 0.004 & low-outgassing FDM \\
SLA resins (Clear, Tough, Grey) & 0.125 (VeroClear) & -- & 1.5--2.9 / $\leq$0.02 & Rigid 10K: 0.54 / 0.00 \\
PVDF / PCTFE        & 0.031 / 0.093 & --       & 0.27 / 0.01 (PVDF) & excellent barriers \\
PTFE / FEP          & 2.3 / 4.5   & --         & -- & poor \Oxy{} barrier despite inertness \\
FKM / butyl O-rings & 1.3--1.7    & 40--200    & 0.10--0.82 / $\leq$0.04 (FKM) & preferred seal materials \\
polyurethane        & 1.5--4.7    & 340--12\,500 & -- & avoid for gas lines \\
silicone            & 100--600    & 10\,500    & -- & never near the end faces \\
\bottomrule
\end{tabularx}
\end{table}


\subsection{Technical facts that shape the testbed}\label{sec:sota-hardware}
\begin{itemize}
  \item \textbf{Source.} The only single-mode \SI{375}{nm} laser diode in Nichia's line-up is the
        NDU4116 (\SI{70}{mW}), sold by Thorlabs as L375P70MLD~\cite{ThorlabsL375}. It needs a TEC
        mount that supports its pin code (e.g.\ LDM56F). Splitting its free-space beam before any fibre
        makes every input face a test item, and no upstream delivery fibre ages. With \SIrange{30}{60}{\percent}
        coupling efficiency, four to six arms at \SI{3}{mW} guided are possible.
  \item \textbf{Detectors.} Silicon photodiodes lose UV responsivity under exposure~\cite{Korde1987,Siles2024};
        photodiodes with high UV resistance exist~\cite{HamamatsuS12698}. SiC photodiodes respond only weakly
        at \SI{369.5}{nm} because it lies at their band edge.
  \item \textbf{Lenses.} The moulded aspheres commonly used for fibre coupling are made of high-index
        glasses whose UV transmission is reduced. Check them at \SI{369.5}{nm}.
  \item \textbf{Inspection standards.} The core zone of IEC~61300-3-35 has a diameter of
        \SI{25}{\micro m}, about ten times the mode-field diameter~\cite{IEC61300335}. Standard
        pass/fail criteria are too coarse for these fibres. An in-house criterion based on
        $\geq400\times$ imaging of the central \SI{10}{\micro m} is needed.
  \item \textbf{Polishing.} Thorlabs' procedure for FC/APC uses diamond films down to \SI{1}{\micro m}
        and a final \SI{0.01}{\micro m} silica film, and warns that alumina and SiC films undercut fibres
        in zirconia ferrules~\cite{ThorlabsPolish}. Every final film and lubricant also leaves a
        possible residue on the face. The procedure is therefore part of the experimental protocol.
\end{itemize}

\subsection{Contribution of this project}
The project combines four elements that, to our knowledge, have not been brought together:
(i) the step-index pure-silica PM fibre and the wavelength and power actually used in \Ybion{}
experiments; (ii) atmosphere as a controlled variable, with \Oxy, water, organic vapour and particles
separated; (iii) simultaneous, identical arms with continuous logging of throughput, face reflection and
gas composition; and (iv) a low-cost, deployable enclosure qualified against the same measurements.
